% BLOCO P55. Descritivo.
\begin{table}[htbp]\centering
\caption{Decomposing the fall in age at administrative resolution, all applicants (2010--2014)}
\label{tab:P55res}
\footnotesize
\setlength{\tabcolsep}{5pt}
\begin{tabular}{lrrrr}
\toprule
Cohort & N & Age at resol. & Age at appl. & Interval \\
\midrule
2010 & 128,352 & 36.73 & 29.07 & 7.64 \\
2011 & 146,693 & 32.67 & 23.43 & 9.23 \\
2012 & 155,292 & 28.85 & 19.96 & 8.87 \\
2013 & 153,033 & 25.40 & 18.17 & 7.20 \\
2014 & 151,772 & 23.36 & 15.51 & 7.83 \\
\midrule
$\Delta$ 2010--14 &  & -13.37 & -13.56 & +0.19 \\
\% of decline &  & 100 & 101 & -1 \\
\bottomrule
\end{tabular}
\\[3pt]{\footnotesize \textit{Note:} Resolution is the first terminal event for each applicant---referral, enrolment, aging out at 47 months without referral, or recorded withdrawal---so it is defined for the \emph{entire} applicant universe, not only enrollees. Age at resolution $=$ age at application $+$ interval, an identity that holds in means (max residual $0.02$ months, i.e.\ under one day, from calendar-day rounding). The fall in mean age at resolution between the 2010 and 2014 cohorts partitions into an application-timing component ($\Delta$ age at application) and an interval component ($\Delta$ time to resolution). The interval is not a processing time: it runs to whichever terminal event comes first and so pools successful outcomes (referral, enrolment) with unsuccessful ones (aging out, withdrawal). This is the version of the decomposition that \emph{includes the never-served}; compare it to the 93\%/7\% split conditional on enrolment.}
\end{table}
